% Options for packages loaded elsewhere
\PassOptionsToPackage{unicode}{hyperref}
\PassOptionsToPackage{hyphens}{url}
\PassOptionsToPackage{dvipsnames,svgnames,x11names}{xcolor}
%
\documentclass[
  11pt,
]{article}
\usepackage{amsmath,amssymb}
\usepackage{iftex}
\ifPDFTeX
  \usepackage[T1]{fontenc}
  \usepackage[utf8]{inputenc}
  \usepackage{textcomp} % provide euro and other symbols
\else % if luatex or xetex
  \usepackage{unicode-math} % this also loads fontspec
  \defaultfontfeatures{Scale=MatchLowercase}
  \defaultfontfeatures[\rmfamily]{Ligatures=TeX,Scale=1}
\fi
\usepackage{lmodern}
\ifPDFTeX\else
  % xetex/luatex font selection
\fi
% Use upquote if available, for straight quotes in verbatim environments
\IfFileExists{upquote.sty}{\usepackage{upquote}}{}
\IfFileExists{microtype.sty}{% use microtype if available
  \usepackage[]{microtype}
  \UseMicrotypeSet[protrusion]{basicmath} % disable protrusion for tt fonts
}{}
\makeatletter
\@ifundefined{KOMAClassName}{% if non-KOMA class
  \IfFileExists{parskip.sty}{%
    \usepackage{parskip}
  }{% else
    \setlength{\parindent}{0pt}
    \setlength{\parskip}{6pt plus 2pt minus 1pt}}
}{% if KOMA class
  \KOMAoptions{parskip=half}}
\makeatother
\usepackage{xcolor}
\usepackage[margin=1in]{geometry}
\setlength{\emergencystretch}{3em} % prevent overfull lines
\providecommand{\tightlist}{%
  \setlength{\itemsep}{0pt}\setlength{\parskip}{0pt}}
\setcounter{secnumdepth}{-\maxdimen} % remove section numbering
\ifLuaTeX
  \usepackage{selnolig}  % disable illegal ligatures
\fi
\usepackage{bookmark}
\IfFileExists{xurl.sty}{\usepackage{xurl}}{} % add URL line breaks if available
\urlstyle{same}
\hypersetup{
  pdftitle={P15.1: Prescribed-type unnormalized energy},
  colorlinks=true,
  linkcolor={Maroon},
  filecolor={Maroon},
  citecolor={Blue},
  urlcolor={Blue},
  pdfcreator={LaTeX via pandoc}}

\title{P15.1: Prescribed-type unnormalized energy}
\author{}
\date{4 October 2026}

\begin{document}
\maketitle

\section{P15.1: prescribed-type unnormalized energy --- partial
result}\label{p15.1-prescribed-type-unnormalized-energy-partial-result}

Authoritative mathematical argument, independent investigation, 4
October 2026. Original atom remains open; no claim concerning P15.2 or
P15.3.

\section{A prescribed bounded-type energy without a
fibration}\label{a-prescribed-bounded-type-energy-without-a-fibration}

\textbf{Status: a partial result for P15.1.} The theorem below applies
to every connected smooth compact Kähler manifold and every actual
reference potential in every pseudoeffective zero-volume class. Its
domain is the prescribed bounded singularity type. It does not assert a
complete solution of the original atom, an identification with numerical
dimension, or any solution of P15.2 or P15.3. All assertions are
mathematical drafts, not independently certified results.

\subsection{Conventions and the standard foundations
used}\label{conventions-and-the-standard-foundations-used}

Let \(X\) be a connected smooth compact Kähler \(n\)-fold, \(n\geq1\).
Set \[
d^c=\frac{i}{4\pi}(\bar\partial-\partial),\qquad dd^c=\frac{i}{2\pi}\partial\bar\partial.
\] A potential is never identically minus infinity. For a smooth closed
real \((1,1)\)-form \(\theta\), put \(T_u=\theta+dd^cu\). Write
\(u\sim\psi\) when \(\psi-C\leq u\leq\psi+C\) for some finite constant;
differences are interpreted on the common finite locus and may be
assigned any value on the common polar set. Define \[
\mathcal C_\psi=\{u\in\operatorname{PSH}(X,\theta):u\sim\psi\}.
\] The comparison inequalities, not an ideal equivalence, define the
retained singularity data.

The following scoped published foundations are used. Their complete
supplied primary bodies, including defects, are reproduced in the
authoritative argument's source annex.

\begin{enumerate}
\def\labelenumi{\arabic{enumi}.}
\tightlist
\item
  \textbf{Nonpluripolar products.} BEGZ, \emph{Monge--Ampère equations
  in big cohomology classes}, Acta Math. 205 (2010), 199--262,
  arXiv:0812.3674v2, Definition 1.1, Proposition 1.4, Proposition 1.6
  and Theorem 1.8: products of finitely many closed positive
  \((1,1)\)-currents on a compact Kähler manifold are well-defined
  finite closed positive currents, symmetric, multilinear on positive
  sums, and plurifinely local. Their trace measures put no mass on
  pluripolar sets. Smooth factors agree with ordinary multiplication.
  The empty product is 1. These statements apply to actual currents and
  do not assert multilinearity of the class positive product.
\item
  \textbf{Bounded-type mixed masses and monotonicity.} Darvas--Di
  Nezza--Lu, \emph{Monotonicity of nonpluripolar products and complex
  Monge--Ampère equations with prescribed singularity}, Anal. PDE 11
  (2018), 2049--2087, Proposition 2.1: for \(n\) smooth closed real
  \((1,1)\)-forms \(\theta_j\) with pseudoeffective classes, and
  \(u_j,v_j\in\operatorname{PSH}(X,\theta_j)\) of the same bounded
  singularity type in every slot, the two top mixed nonpluripolar
  integrals are equal. Theorem 2.4: if every \(u_j\) is less singular
  than \(v_j\), the integral for the \(u_j\)'s is at least that for the
  \(v_j\)'s. There is no big-class or small-unbounded-locus hypothesis
  in these statements. Smooth Kähler factors are slots with potential
  zero.
\item
  \textbf{Volume at the boundary.} BEGZ Definition 1.17 defines the
  volume of a big class by its minimal-singularity top mass and extends
  the positive product to pseudoeffective classes as the decreasing
  limit along addition of an arbitrary Kähler class. In particular
  \(\operatorname{vol}(\alpha+\varepsilon\lambda)\to\operatorname{vol}(\alpha)\)
  as \(\varepsilon\downarrow0\). For nef classes the positive product
  equals the ordinary cup product. To bound the mass of an arbitrary
  current in a big class we use foundation 2: a minimal-singularity
  potential is less singular than it, so its top mass is an upper bound.
  No equality with the arbitrary actual-current mass is assumed.
\item
  \textbf{Admissible Dirichlet products and IBP.} Duc-Viet Vu,
  \emph{Convexity of the class of currents with finite relative energy},
  arXiv:2005.13241v3 (25 March 2022), §2, Lemmas 2.5--2.6 and Theorem
  2.7: if \(R\) is a closed positive \((n-1,n-1)\)-current and (a,b) are
  bounded \(R\)-admissible dsh functions, then their intrinsic relative
  Dirichlet products exist; the diagonal product is positive, and \[
  \int b\langle dd^ca\dot\wedge R\rangle
  =\int a\langle dd^cb\dot\wedge R\rangle
  =-\int\langle db\wedge d^ca\dot\wedge R\rangle.
  \tag{F}
  \] A dsh function is a difference of two quasi-psh functions.
  Admissibility means that there is such a representation whose polar
  sets have zero \(R\)-mass. Vu's Proposition 2.1(iii) identifies
  iterated relative products; with
  \(R=\langle S_1\cdots S_m\rangle\wedge W\), \(W\) smooth closed
  positive, the signed product in (F) is the difference of the
  corresponding ordinary nonpluripolar products. Products in all our
  applications are finite by foundation 1. If (m=0), (R=W) is smooth.
  Otherwise its trace measure has no pluripolar mass. Thus every bounded
  difference used below is admissible, including smooth test functions.
\end{enumerate}

Vu's current convention is \(dd^c_V=i\partial\bar\partial/\pi=2dd^c\);
the read-only page-4 image of arXiv:2004.11111v1 confirms it despite a
defective text extraction. A source potential \(u/2\) represents the
current \(T_u\), with the smooth background unchanged. Source
differences are divided by 2; \(d^c_V=2d^c\). The signed source
\(dd^c_V(a/2)\) equals \(dd^ca\), and the source Dirichlet current for
(a/2) equals one half of ours. Multiplying the source IBP equality by 2
gives (F). Source bounded-type comparisons, full-mass conditions, and
energy finiteness are unchanged by this positive scalar change. DDNL
uses \(i\partial\bar\partial=2\pi dd^c\); its source potentials are
\(u/(2\pi)\). Its mixed-mass theorems concern the same actual currents
after that conversion. The locally psh BEGZ product statements are
likewise independent of a positive convention change, by rescaling local
potentials to represent the same currents. No application of Xia's
big-class corollary to a merely pseudoeffective class is made here.

The intrinsic diagonal definition is \[
da\wedge d^ca\wedge R=\tfrac12dd^c(a^2)\wedge R-a\,dd^ca\wedge R
\tag{G}
\] for differences of bounded quasi-psh functions, followed by Vu's
admissible nonpluripolar truncation construction. The factor (1/2) is
mandatory. Vu v3 Eq. (2.3) prints it correctly, but the repeated page-6
expansion and page-11 cutoff identity omit it. We preserve those literal
quotations and use (G), not those defective repetitions.

For clarity, the defect does not introduce a gap in our use of (F). Here
is a direct corrected cutoff check. For negative quasi-psh
representations \(a=\phi_1-\phi_2\), \(b=\phi_3-\phi_4\), choose a
Kähler form \(\Omega\) with \(dd^c\phi_j\geq-\Omega\). Put \[
\chi_k=1+k^{-1}\max(\phi_1+\phi_2+\phi_3+\phi_4,-k),\quad
a_k=\max(\phi_1,-k)-\max(\phi_2,-k),
\] and define \(b_k\) similarly. Then \(0\leq\chi_k\leq1\),
\(dd^c\chi_k\geq-4\Omega/k\), and the support of \(\chi_k\) is in the
simultaneous finite truncation locus. Correct bounded
Bedford--Taylor/Stokes calculus using (G) gives \[
0\leq\int d\chi_k\wedge d^c\chi_k\wedge R
=-\int\chi_kdd^c\chi_k\wedge R
\leq\frac4k\int\Omega\wedge R.
\tag{C}
\] Indeed \(dd^c\chi_k+4\Omega/k\) is positive, so its negative
contribution in the last expression can be discarded. The same
calculation gives
\(\int da_k\wedge d^ca_k\wedge R\leq2\|a_k\|_\infty\int\Omega R\); the
sup norms are uniformly bounded by boundedness of \(a\). These formulae
apply to bounded quasi-psh truncations by bounded Bedford--Taylor
calculus. Cauchy--Schwarz and (C) make all cross terms containing
\(d\chi_k\) tend to zero. On \(\{\chi_k>0\}\), plurifine locality
identifies the signed products and the intrinsic gradients with their
truncations. Dominated convergence is permissible because \(R\) has no
mass on the representation poles, the relative positive products also
have no mass there, and the admissible gradients have finite mass by
Vu's Lemmas 2.5--2.6. Applying bounded IBP to
\(\chi_k(b_kdd^ca_k-a_kdd^cb_k)R\) proves symmetry; applying it to
\(\chi_k(a_kdd^ca_k+da_k\wedge d^ca_k)R\) proves the negative diagonal
equality. Polarization gives (F). In the page-6 local square expansion,
the correct coefficients of its three square terms are \(1,1,-1/2\), not
\(2,2,-1\). Thus every Dirichlet calculation here has the correct
intrinsic factor. This check uses the standard gradient construction,
not the misprinted formula as a definition.

\subsection{The bounded theorem}\label{the-bounded-theorem}

Suppose \(\alpha=[\theta]\) is pseudoeffective with
\(\operatorname{vol}(\alpha)=0\), and fix an actual
\(\psi\in\operatorname{PSH}(X,\theta)\). For a Kähler form \(\omega\),
set \[
M_k(\psi,[\omega])=\int_X\langle T_\psi^k\rangle\wedge\omega^{n-k},\quad
d=\max\{0\leq k\leq n:M_k>0\},\quad r=n-d.
\tag{1}
\] The convention is \(M_0=\int\omega^n>0\). Then \(d<n\), \(d\) is
independent of the auxiliary Kähler form, and every degree
(k\textgreater d) nonpluripolar mixed product of currents \(T_{u_i}\),
\(u_i\in\mathcal C_\psi\), is zero as a current.

For any Kähler class \(\lambda\), the following relative energy is
independent of the Kähler representative \(\omega\in\lambda\): \[
\mathcal E_{\psi,\lambda}(u)
=\frac1{d+1}\sum_{a=0}^d\int_X(u-\psi)
\langle T_u^aT_\psi^{d-a}\rangle\wedge\omega^r.
\tag{2}
\] For any \(u,v\in\mathcal C_\psi\), let \(\mathcal E_{v,\lambda}(u)\)
denote the same formula with anchor \(v\). Then \[
\mathcal E_{\psi,\lambda}(u)-\mathcal E_{\psi,\lambda}(v)
=\mathcal E_{v,\lambda}(u).
\tag{3}
\] The energy is increasing, concave,
\(\mathcal E_{\psi,\lambda}(\psi)=0\), and \[
\mathcal E_{\psi,\lambda}(u+c)=\mathcal E_{\psi,\lambda}(u)+cM_d,
\qquad
|\mathcal E_{\psi,\lambda}(u)-\mathcal E_{\psi,\lambda}(v)|
\leq M_d\|u-v\|_\infty.
\tag{4}
\] No division by zero volume, or division by \(M_d\), occurs. The
carrier for these integrals is the actual \(n\)-dimensional manifold
\(X\); the nonpluripolar current has bidegree \((d,d)\). A geometric
\(d\)-dimensional quotient is not supplied or assumed.

\subsubsection{Degree, mass, and
vanishing}\label{degree-mass-and-vanishing}

The class \(\alpha+\varepsilon[\omega]\) is big for \(\varepsilon>0\):
the actual current \(T_\psi+\varepsilon\omega\) is bounded below by
\(\varepsilon\omega\). A minimal-singularity potential in this big class
is less singular than \(\psi\), so foundation 2 gives \[
0\leq M_n\leq\int\langle(T_\psi+\varepsilon\omega)^n\rangle
\leq\operatorname{vol}(\alpha+\varepsilon[\omega]).
\] The middle expression expands, using actual-current multilinearity,
as \[
P_\varepsilon(\omega)=\sum_{k=0}^n\binom nk\varepsilon^{n-k}M_k.
\tag{5}
\] The last bound tends to zero by foundation 3. Hence \(M_n=0\),
\(d<n\), and (5) ends at \(k=d\). This does not identify
\(P_\varepsilon\) with class volume.

If \(\omega'\) is another Kähler form, compactness gives constants
\(c,C>0\) such that \(c\omega\leq\omega'\leq C\omega\). For a positive
\((k,k)\)-current \(S\), positivity of the power differences gives \[
c^{n-k}\int S\omega^{n-k}\leq\int S(\omega')^{n-k}\leq C^{n-k}\int S\omega^{n-k}.
\] Thus the positivity tests in (1) agree. Foundation 2, with \(k\)
same-type current slots and (n-k) slots \((\omega,0)\), gives \[
\int\langle T_{u_1}\cdots T_{u_k}\rangle\omega^{n-k}=M_k.
\tag{6}
\] For \(k>d\) the left current is positive and has zero trace mass
against a strictly positive form. A positive current of zero trace mass
is zero: on each coordinate chart its coefficient measures are dominated
by its trace measure, by positivity. This proves the stated current
vanishing. Each \(M_k\) depends only on \([\omega]\) because
\(\langle T_\psi^k\rangle\) is closed and integration of an exact smooth
complementary form against a closed current vanishes by distributional
Stokes. It is a homogeneous polynomial of degree (n-k) in \([\omega]\).

\subsubsection{Closed polynomial one-form and its
primitive}\label{closed-polynomial-one-form-and-its-primitive}

We first work at any degree \(0\leq q\leq n\), with a fixed smooth
product \(W=\omega^{n-q}\). In this paragraph there is no zero-volume
assumption. On the convex hull of a finite family \(u_0,\ldots,u_m\) of
same-type \(\theta\)-psh potentials, set \(h_i=u_i-u_0\),
\(u_t=u_0+\sum t_i h_i\), and \[
A_i(t)=\int h_i\langle T_{u_t}^q\rangle W.
\] Multilinearity makes these functions polynomials in the finite
parameters. All coefficient integrals exist since \(h_i\) is bounded off
a common polar set and the measures do not charge it. For \(q\geq1\), \[
\partial_j A_i=q\int h_i\langle dd^ch_j\wedge T_{u_t}^{q-1}\rangle W
=\partial_i A_j
\tag{7}
\] by (F), applied to the closed positive current
\(R=\langle T_{u_t}^{q-1}\rangle W\). Its bidegree is \((n-1,n-1)\); it
has no pluripolar mass, including when \(q=1\), in which case it is
smooth. The signed expression is exactly the difference in one current
slot, by the iterated-product property. Thus (7) is a checked
application, not formal unsigned wedge calculus. For \(q=0\), the
\(A_i\)'s are constant and closedness is immediate.

The line integral of this closed one-form from \(v\) to \(u\) equals \[
\int_0^1\int (u-v)\langle((1-t)T_v+tT_u)^q\rangle W\,dt
=\frac1{q+1}\sum_{a=0}^q\int(u-v)\langle T_u^aT_v^{q-a}\rangle W.
\tag{8}
\] The last equality follows by the binomial expansion and
\(\binom qa\int_0^1t^a(1-t)^{q-a}dt=1/(q+1)\). Closedness (7) implies
zero integral around each triangle in the convex hull, by ordinary
Stokes for a polynomial one-form on that finite simplex. The identity
extends to its boundary by polynomial continuity. Taking the triangle
with vertices \(\psi,v,u\) proves (3), for degree \(d\), and proves the
corresponding cocycle at every other degree or positive perturbation
used later.

For a segment \(u_t=(1-t)u+tv\), \(h=v-u\), this also proves exact
variation regularity: \[
\frac d{dt}\mathcal E_{\psi,\lambda}(u_t)=\int h\langle T_{u_t}^d\rangle\omega^r,
\tag{9}
\] and, for \(d\geq1\), \[
\frac {d^2}{dt^2}\mathcal E_{\psi,\lambda}(u_t)
=d\int h\langle dd^ch\wedge T_{u_t}^{d-1}\rangle\omega^r
=-d\int\langle dh\wedge d^ch\dot\wedge
(\langle T_{u_t}^{d-1}\rangle\omega^r)\rangle\leq0.
\tag{10}
\] The last inequality is positivity of the intrinsic Dirichlet product
(G); IBP symmetry alone would not give it. Equations (9)--(10) are
polynomial identities with continuous endpoint derivatives. More
generally the pullback to any finite simplex is polynomial and admits
the ordinary chain rule along \(C^1\) parameter paths and second
derivatives along \(C^2\) parameter paths. No derivative assertion for
arbitrary discontinuous or unbounded velocities is made. When \(d=0\),
the energy is affine and there is no gradient formula (10).

If \(u\leq v\), (9) is nonnegative. A constant shift leaves the current
fixed, and (6) makes (9) equal \(cM_d\). Finally, the pointwise
sandwiches \(v-\delta\leq u\leq v+\delta\), \(\delta=\|u-v\|_\infty\),
give (4). The same argument proves the corresponding Lipschitz estimate
with the actual common-type mass in an arbitrary positive perturbation.

\subsubsection{Representative independence and cohomological
continuity}\label{representative-independence-and-cohomological-continuity}

If \(\omega_1,\omega_0\) are cohomologous Kähler forms, the compact
Kähler \(dd^c\)-lemma gives a smooth real \(g\) with
\(\omega_1-\omega_0=dd^cg\). We use exactly the degree-two consequence
of this standard lemma: cohomologous smooth closed real \((1,1)\)-forms
differ by \(dd^c\) of a smooth real function. A standard published
reference is Deligne--Griffiths--Morgan--Sullivan, \emph{Real homotopy
theory of Kähler manifolds}, Invent. Math. 29 (1975), 245--274, §5, the
principle of two types. The hypotheses here are precisely compact
Kähler, smooth, closed, real, cohomologous, and of pure type \((1,1)\);
taking the real part and rescaling the operator gives our stated
convention. This is a standard foundation citation, not an independently
reproduced source proof. For \(r\geq1\), \[
\omega_1^r-\omega_0^r=dd^cg\wedge\sum_{b=0}^{r-1}\omega_1^b\omega_0^{r-1-b}.
\] For any term in (2), write \(h=u-\psi\) and
\(Q=\langle T_u^aT_\psi^{d-a}\rangle\). For each \(b\), the current
\(R=Q\omega_1^b\omega_0^{r-1-b}\) is closed positive of bidegree
\((n-1,n-1)\). It has no pluripolar mass; when \(d=0\), it is smooth.
Thus \(h,g\) are bounded admissible dsh functions and (F) gives \[
\int h\,dd^cg\,R=\int g\langle dd^ch\dot\wedge R\rangle=0.
\tag{11}
\] The signed current on the right is
\((\langle T_uT_u^aT_\psi^{d-a}\rangle-
\langle T_\psi T_u^aT_\psi^{d-a}\rangle)\omega_1^b\omega_0^{r-1-b}\);
both terms vanish by degree \(d+1\) vanishing. All multiplication by
bounded functions takes place against nonpluripolar measures. This
proves representative independence, including the \(d=0\) case.

The same proof applies to a mixed product of \(r\) smooth closed Kähler
factors, changing any one factor within its class. It makes the energy a
symmetric multilinear functional of their classes. To see its polynomial
dependence without imposing convergence of representatives, choose a
basis \(e_i\) of \(H^{1,1}(X,\mathbb R)\) and fixed smooth closed
representatives \(\eta_i\). Near any fixed Kähler class \(\lambda_0\)
represented by a positive \(\omega_0\), the linear choice
\(\omega(\lambda)=\omega_0+\sum x_i\eta_i\) remains positive for small
\(x_i\). The expression (2) is a polynomial in these coordinates.
Alternatively decompose arbitrary smooth closed real factors as
differences of Kähler forms by adding a large fixed Kähler form, and
apply multilinearity. Representative independence ensures that these
local polynomials agree on overlap and give one homogeneous polynomial
of degree \(r\) on the Kähler cone. In particular the leading energy is
continuous in the auxiliary class.

Changing \(\theta\) to \(\theta+dd^cg\) and simultaneously replacing
\(u,\psi\) by \(u-g,\psi-g\) leaves every actual current and difference
unchanged. Changing the anchor within \(\mathcal C_\psi\) has exactly
the cocycle (3), not absolute anchor independence.

\subsection{Relative Kähler perturbation
limits}\label{relative-kuxe4hler-perturbation-limits}

For \(\varepsilon>0\), an arbitrary Kähler form \(\omega\), and
\(v\in\operatorname{PSH}(X,\theta+\varepsilon\omega)\) with
\(v\sim\psi\), define the unnormalized relative top energy \[
\mathcal E^{n,\varepsilon,\omega}_\psi(v)
=\frac1{n+1}\sum_{a=0}^n\int(v-\psi)
\langle(T_v+\varepsilon\omega)^a(T_\psi+\varepsilon\omega)^{n-a}\rangle.
\tag{12}
\] Here \(T_v\) alone need not be positive; the displayed currents
\(T_v+\varepsilon\omega\) are positive. The degree-\(n\) version of
(7)--(8), applied directly in the pseudoeffective (indeed big) class
\(\alpha+\varepsilon[\omega]\), gives its cocycle, monotonicity and
Lipschitz constant \(P_\varepsilon(\omega)\), the actual mass (5). It is
not class volume in general.

For fixed \(u\in\mathcal C_\psi\), put \[
\mathcal E^{(k),\omega}_\psi(u)=\frac1{k+1}\sum_{a=0}^k
\int(u-\psi)\langle T_u^aT_\psi^{k-a}\rangle\omega^{n-k}.
\] Expanding the segment formula (8) gives the exact identity \[
\mathcal E^{n,\varepsilon,\omega}_\psi(u)
=\sum_{k=0}^d\binom nk\varepsilon^{n-k}\mathcal E^{(k),\omega}_\psi(u).
\tag{13}
\] In the original mixed sum, the coefficient for
\(T_u^aT_\psi^{k-a}\omega^{n-k}\) is
\(\sum_{j=0}^n\binom ja\binom{n-j}{k-a}/(n+1)
=\binom n k/(k+1)\). The identity follows by marking the \((a+1)\)-st
element of a \((k+1)\)-subset of \({0,\ldots,n}\) and summing over its
location. This also proves it at \(k=0\). For all \(k\), (6) gives \[
|\mathcal E^{(k),\omega}_\psi(u)|\leq\|u-\psi\|_\infty M_k(\psi,[\omega]).
\tag{14}
\]

Now let \textbf{any} family of Kähler forms \(\omega_\varepsilon\) have
\([\omega_\varepsilon]\to\lambda\), with \(\lambda\) Kähler. No smooth
or pointwise convergence of these forms is assumed. Let
\(u_\varepsilon\in\operatorname{PSH}(X,\theta+\varepsilon\omega_\varepsilon)\)
have the same bounded type as \(\psi\), and let
\(\delta_\varepsilon=\|u_\varepsilon-u\|_\infty\to0\) off their common
polar set. Then \[
\frac{\mathcal E^{n,\varepsilon,\omega_\varepsilon}_\psi(u_\varepsilon)}
{\binom n d\varepsilon^{n-d}}
\longrightarrow\mathcal E_{\psi,\lambda}(u).
\tag{15}
\] Indeed the positive-\(\varepsilon\) cocycle and monotonicity give \[
|\mathcal E^{n,\varepsilon,\omega_\varepsilon}_\psi(u_\varepsilon)
-\mathcal E^{n,\varepsilon,\omega_\varepsilon}_\psi(u)|
\leq\delta_\varepsilon P_\varepsilon(\omega_\varepsilon).
\] Each \(M_k\) is a cohomological polynomial and is bounded for classes
in a compact neighborhood of \(\lambda\); hence
\(P_\varepsilon=O(\varepsilon^r)\), uniformly without any bounds on
representatives. The normalized error tends to zero. For fixed \(u\),
(13)--(14) bound all normalized \(k<d\) terms by constants times
\(\varepsilon^{d-k}\). The \(k=d\) term is
\(\mathcal E_{\psi,[\omega_\varepsilon]}(u)\), converging by its class
polynomial. This proves (15). Only the leading coefficient is
representative independent. Formula (13) does not claim representative
independence of the lower coefficients or of the full energy (12).

For example, on \(X=\mathbb P^1\times\mathbb P^1\), take
\(\theta=p_1^*\kappa\), \(\psi=0\), and a nonconstant smooth
\(u=p_1^*v\) with \(\kappa+dd^cv>0\). Set
\(\omega_0=p_1^*\kappa+p_2^*\kappa\), with \(\int\kappa=1\). Replacing
\(\omega_0\) by \(\omega_0+tdd^cu\), for small nonzero \(t\), stays
Kähler and in the same class. Here \(d=1\); the leading
\(\mathcal E^{(1)}\) is unchanged, but \[
\int u\big((\omega_0+tdd^cu)^2-\omega_0^2\big)
=2t\int_{\mathbb P^1}vdd^cv
=-2t\int_{\mathbb P^1}dv\wedge d^cv\ne0.
\] Thus the \(\varepsilon^2\) coefficient \(\mathcal E^{(0)}\) changes.
This explicitly rules out an unintended claim about the full
positive-perturbation energy.

\subsection{Auxiliary differential-measure check, only for the energy
route}\label{auxiliary-differential-measure-check-only-for-the-energy-route}

No zero-volume operator theorem is claimed here. The following estimate
checks the differential in (9) under a \textbf{fixed} Kähler form. For
fixed \(u\in\mathcal C_\psi\), (5) as a measure expansion and degree
vanishing give \[
\frac{\langle(T_u+\varepsilon\omega)^n\rangle}
{\binom n d\varepsilon^r}
\rightharpoonup\langle T_u^d\rangle\omega^r.
\tag{16}
\] The limiting measure is on \(X\), has positive raw mass \(M_d\), and
still depends on the full form in general.

For \(u_\varepsilon\) as in (12), set
\(A=T_{u_\varepsilon}+\varepsilon\omega\), \(B=T_u+\varepsilon\omega\),
\(h=u_\varepsilon-u\), \(\delta=\|h\|_\infty\). Multilinearity factors
the top difference as
\(\sum_{j=0}^{n-1}\langle dd^ch\wedge A^jB^{n-1-j}\rangle\). For a
smooth test \(f\), (F) transfers \(dd^c\) to \(f\), with admissible
bounded \(h\), and \(-C_f\omega\leq dd^cf\leq C_f\omega\) bounds the
result by \[
\left|\int f(\langle A^n\rangle-\langle B^n\rangle)\right|
\leq nC_f\delta\,Q_\varepsilon,
\quad
Q_\varepsilon=\sum_{k=0}^d\binom{n-1}k\varepsilon^{n-1-k}M_k
=O(\varepsilon^{r-1}).
\tag{17}
\] For every \(j\), equality of the mixed mass with one smooth
\(\omega\)-slot to \(Q_\varepsilon\) is exactly foundation 2 in class
\(\alpha+\varepsilon[\omega]\), replacing the \(A,B\)-potentials by
\(\psi\), followed by positive binomial expansion. Since \(d<n\), the
sum has legal indices. At \(n=1,d=0\), the product with zero
\(A,B\)-slots is 1 and \(Q_\varepsilon=M_0\), so the same calculation
applies. The normalized error is \(O_f(\delta/\varepsilon)\). Therefore
\(\delta=o(\varepsilon)\) suffices for (16) with moving potentials; this
is a sufficient rate, not an optimality claim.

If \(\omega_\varepsilon\to\omega\) smoothly and \(\omega\) is Kähler,
the tests are uniformly bounded by these forms and the same constants
can be taken uniformly; the fixed-current terms in (16) converge by
smooth multiplication. For class-convergent forms alone, neither this
test bound nor convergence of the final full-form measure is supplied.
The rate-free energy result (15) is not transferred to (16).

\section{A genuine relative finite-energy extension, with specified
continuity}\label{a-genuine-relative-finite-energy-extension-with-specified-continuity}

This note uses the bounded theorem proved in local note \textbf{n1}. It
retains the same actual reference type and auxiliary Kähler class. It
proves a finite-energy extension of that theorem, with decreasing
monotone continuity. It does not claim an unbounded perturbative limit
or arbitrary monotone/path regularity.

\subsection{The precise additional standard
foundations}\label{the-precise-additional-standard-foundations}

Vu, arXiv:2005.13241v3, §3, with \(P=T_\psi\), \(m=d\),
\(T=\omega^{n-d}\), and linear weight \(\chi(t)=t\), defines the
following spaces. A current \(S\) in \([P]\), more singular than \(P\),
has relative full mass if \[
\{\langle S^d\dot\wedge T\rangle\}=
\{\langle P^d\dot\wedge T\rangle\}.
\] For negative potentials \(u\leq\psi\), its diagonal joint
linear-weight energy is \[
d\sum_{a=0}^d\binom da\int(\psi-u)
\langle T_u^aT_\psi^{d-a}\rangle\omega^{n-d}.
\tag{16a}
\] This is the exact specialization of Vu's Eq. (3.1): there are
\(\binom da\) subsets of size \(a\), and the weight is
\(\chi(d(u-\psi))=d(u-\psi)\). Theorem 3.4 states that both the
full-mass space \(\mathcal E_{d,P}(T)\) and its joint weighted subclass
\(\mathcal E_{\chi,d,P}(T)\) are convex. Corollary 3.2 states that the
weighted class is unchanged when \(P\) is replaced by a cohomologous
current of the same bounded singularity type. Finiteness does not depend
on additive potential constants. Here \(T\) is smooth closed positive,
the relative products are ordinary nonpluripolar products with smooth
factors, and the product has top degree. On connected \(X\), equality of
its cohomology classes is equality of total masses. Thus all source
hypotheses are satisfied, with no minimal-singularity or model-envelope
assumption. The source proofs of convexity and reference-type invariance
use cited arguments; they are used as scoped standard foundations, not
as reconstructed new lemmas. The convention conversion in n1 preserves
full mass and weighted finiteness, multiplying (16a) by a positive
constant.

We also use the following standard bounded convergence foundations
solely to prove the continuity assertion below. A decreasing sequence of
quasi-psh functions in a fixed smooth background, with a quasi-psh limit
not identically minus infinity, converges in Kähler Monge--Ampère
capacity. This is the decreasing-convergence statement cited as
Guedj--Zeriahi 2005, Proposition 3.7, in DDNL's proof of Theorem 2.4,
p.~2057. DDNL Proposition 2.2, pp.~2054--2055, states: on an open set in
\(\mathbb C^n\), uniformly bounded local psh functions converging in
capacity in every slot, multiplied by uniformly bounded quasicontinuous
functions also converging in capacity, give weakly convergent weighted
Bedford--Taylor top products. Quasi-psh functions and their bounded
truncations are quasicontinuous. All applications below have a fixed
background and uniformly bounded \textbf{local truncations}. We do not
apply the big-class-only DDNL Theorem 2.3 directly to the
pseudoeffective background. The complete primary body of Proposition 2.2
is included in the argument annex; the scalar capacity fact is a
standard foundation cited through the supplied DDNL proof, not a new
source-proven theorem here.

\subsection{Domain and theorem}\label{domain-and-theorem}

Fix \(\omega\in\lambda\), \(r=n-d\), and \(M=M_d(\psi,\lambda)>0\). For
\(d\geq1\), let \(\mathcal F^1_{\psi,\lambda}\) consist of
\(\theta\)-psh \(u\) such that \(u\leq\psi+C\), \[
\int\langle T_u^d\rangle\omega^r=M,
\qquad
\sum_{a=0}^d\binom da\int|u-\psi|
\langle T_u^aT_\psi^{d-a}\rangle\omega^r<\infty.
\tag{17a}
\] For \(d=0\), use the same more-singular condition and
\(\int|u-\psi|\omega^n<\infty\); degree-zero full mass is automatic.
This integral is finite for quasi-psh \(u,\psi\) since both are locally
integrable against smooth volume. Every bounded-type potential belongs
to the domain.

Define, for this domain, \[
\mathcal E_{\psi,\lambda}(u)
=\frac1{d+1}\sum_{a=0}^d\int(u-\psi)
\langle T_u^aT_\psi^{d-a}\rangle\omega^r.
\tag{18a}
\] Then all these integrals are finite. The energy equals the decreasing
limit of bounded-type truncation energies \(u_L=\max(u,\psi-L)\), is
independent of the truncation levels tending to infinity, and is
continuous under \textbf{every decreasing sequence in
\(\mathcal F^1_{\psi,\lambda}\) with limit in that domain}. In
particular, every decreasing sequence of bounded-type potentials
converging to a point of this domain gives the same extension value. It
is increasing and concave. Its constant shifts are \(cM\). For any two
members \(u,v\) of the domain, all integrals \[
\mathcal E_{v,\lambda}(u)
=\frac1{d+1}\sum_{a=0}^d\int(u-v)
\langle T_u^aT_v^{d-a}\rangle\omega^r
\tag{19a}
\] are finite and obey the exact cocycle \[
\mathcal E_{\psi,\lambda}(u)-\mathcal E_{\psi,\lambda}(v)
=\mathcal E_{v,\lambda}(u).
\tag{20a}
\] Changing the anchor within the original bounded type leaves the
finite-energy domain unchanged and changes the energy by this cocycle.
Both domain and energy are independent of the Kähler representative of
the fixed class \(\lambda\). No identity of unbounded first or second
variations is asserted.

\subsection{Mixed full mass and mixed weight
integrability}\label{mixed-full-mass-and-mixed-weight-integrability}

First subtract constants from finitely many potentials so that all are
\(\leq\psi\), and shift \(\psi\) to be negative if needed. This only
translates energy by constants times finite masses, and will be reversed
at the end. For \(d\geq1\), the definition (17a) is exactly the source
joint linear-weight class (16a), so its convexity applies.

Given \(u_1,\ldots,u_N\) in this domain, including any copies of
\(\psi\), put \(w=N^{-1}\sum_i u_i\). Convexity gives
\(\int\langle T_w^d\rangle\omega^r=M\). By expansion it is the mean of
\(N^d\) ordered mixed masses. Each is at most \(M\), by DDNL Theorem 2.4
replacing every more-singular \(u_i\) by \(\psi\), with \(r\) smooth
Kähler slots. They are nonnegative and their average is \(M\). Hence
\textbf{every} such mixed mass equals \(M\). This proves all mixed full
masses used below without inferring them from the diagonal by an
unsupported polarization.

The \(a=d\) term of the finite joint energy of \(w\) shows that
\(\int(\psi-w)\langle T_w^d\rangle\omega^r<\infty\). Every summand in
its positive expansion is therefore finite. Since \[
0\leq\psi-u_i\leq N(\psi-w),
\] each \(\psi-u_i\) is integrable against every degree-\(d\) mixed
product formed from the finite family. In particular
\(|u-v|\leq(\psi-u)+(\psi-v)\) is integrable against every product in
(19a). All representation poles are pluripolar and have zero mass. For
\(d=0\), these assertions follow directly from smooth-volume
integrability and the common mass \(M_0\).

For completeness, every degree \(k>d\) mixed product of such
more-singular potentials also vanishes: DDNL mass monotonicity bounds
its trace integral by the corresponding zero \(M_k(\psi,\lambda)\), and
positivity implies vanishing as in n1. This assertion does not give
weighted control of any degree \(k<d\).

\subsection{The exact truncation
identity}\label{the-exact-truncation-identity}

Assume \(u\leq\psi\), write \(z=\psi-u\geq0\), and let \[
\mu_a=\langle T_u^aT_\psi^{d-a}\rangle\omega^r,
\qquad
\mu_{a,L}=\langle T_{u_L}^aT_\psi^{d-a}\rangle\omega^r.
\] All these measures have mass \(M\), by the preceding argument and
bounded-type equality respectively. On the plurifine comparison set
\(G_L=\{u>\psi-L\}\), locality gives \(\mu_{a,L}=\mu_a\). The
comparison-set version with unbounded psh data is also exactly Vu v3
Theorem 3.3, applied in local charts to the potentials plus their smooth
local backgrounds. On the complement \(B_L\), off the polar sets,
\(u_L-\psi=-L\). Equality of total masses gives
\(\mu_{a,L}(B_L)=\mu_a(B_L)\). Consequently \[
\int(\psi-u_L)d\mu_{a,L}
=\int_{G_L}z\,d\mu_a+L\mu_a(B_L)
=\int\min(z,L)d\mu_a.
\tag{21a}
\] For \(a=0\), the measure is fixed and the identity is immediate.
Boundary sets where \(u=\psi-L\) cause no issue: the weight is exactly
\(L\) there, and locality is used only on the strict comparison set.
Measures give no mass to infinite-value sets. Monotone convergence in
the right side proves \[
\mathcal E_{\psi,\lambda}(u_L)\downarrow\mathcal E_{\psi,\lambda}(u).
\tag{22a}
\] This proves existence and independence for all real truncation levels
tending to infinity. Order of potentials is inherited by their
truncations, so bounded monotonicity and (22a) prove monotonicity on the
whole domain. The exact mixed masses prove the \(cM\) constant-shift law
directly in (18a); normalization by constants in the proof is therefore
harmless.

\subsection{A bounded decreasing continuity lemma and its
proof}\label{a-bounded-decreasing-continuity-lemma-and-its-proof}

Suppose \(v_j\downarrow v\), all have the bounded type of \(\psi\), and
\(|v_j-\psi|,|v-\psi|\leq B\) uniformly on their finite loci. For a
fixed degree \(q\), each mixed measure formed from \(v_j\) and \(\psi\),
with \(n-q\) fixed smooth Kähler factors, has the same mass as its
limit, by bounded-type equality. Let \(\eta_L=\min(1,\max(\psi+L,0))\).
Then \(0\leq\eta_L\leq1\), it vanishes on \(\{\psi\leq-L\}\), and tends
to 1 outside the polar set of \(\psi\).

On a coordinate chart with smooth local potential \(\rho\) for
\(\theta\), the functions \(\rho+v_j\), \(\rho+v\), and \(\rho+\psi\)
are psh. Where \(\eta_L\ne0\), they have a uniform lower bound depending
only on \(L,B\) and the compact chart. Replace them by local bounded psh
truncations below a still lower constant; locality preserves all
products on that set. Their upper bounds are uniform from compactness
and the fixed \(v_1\). Decreasing quasi-psh convergence gives
convergence in capacity, and the local bounded truncations and their
differences converge in capacity as well. The cutoff and the bounded
weights can be expressed using these truncations on their support, so
they are uniformly bounded quasicontinuous functions converging in
capacity. DDNL Proposition 2.2 with the remaining fixed smooth slots and
a smooth partition of unity therefore proves convergence of both
\(\int\eta_Ld\nu_j\) and \(\int\eta_L(v_j-\psi)d\nu_j\) to their
limiting counterparts. This is a local bounded application in every
slot; no big-class convergence result is needed.

If the common mass is \(m\), then
\(\int(1-\eta_L)d\nu_j=m-\int\eta_Ld\nu_j\). The integrals of the
omitted weights have absolute value at most \(B\) times this quantity.
First let \(j\to\infty\), and then \(L\to\infty\); the limiting tail
mass tends to zero because the limit has no polar mass. This proves
convergence of each full weighted mixed integral and hence convergence
of the bounded-type energy. The same proof allows a fixed same-type
anchor in place of \(\psi\). It explicitly controls the unbounded
reference's polar-locus tails.

\subsection{Decreasing continuity and concavity of the
extension}\label{decreasing-continuity-and-concavity-of-the-extension}

Let \(v_j\downarrow u\) with all \(v_j,u\) in the finite-energy domain.
The sequence has a common upper bound \(\psi+C\), since
\(v_1\leq\psi+C\). By monotonicity, \[
\mathcal E(u)\leq\mathcal E(v_j)\leq\mathcal E(\max(v_j,\psi-L)).
\] For fixed \(L\), the potentials on the right decrease to
\(\max(u,\psi-L)\) within a uniform bounded type. The lemma gives
\(\limsup_j\mathcal E(v_j)\leq\mathcal E(u_L)\). Letting \(L\to\infty\)
in (22a) proves \(\mathcal E(v_j)\to\mathcal E(u)\). This also proves
independence from arbitrary decreasing bounded-type approximants, rather
than merely renamed truncations.

For \(u,v\) in the domain and \(t\in[0,1]\), Vu's convexity makes
\(w=(1-t)u+tv\) a domain member. The bounded-type combinations
\(w_L=(1-t)u_L+tv_L\) decrease to \(w\). Bounded concavity gives
\(\mathcal E(w_L)\geq(1-t)\mathcal E(u_L)+t\mathcal E(v_L)\). Decreasing
continuity and (22a) prove concavity of the extension. For \(d=0\) the
formula is directly affine.

\subsection{Cocycle and reference
changes}\label{cocycle-and-reference-changes}

Take normalized \(u,v\leq\psi\) and their simultaneous truncations at
level \(L\). Let \(\nu_a=\langle T_u^aT_v^{d-a}\rangle\omega^r\),
\(\nu_{a,L}=\langle T_{u_L}^aT_{v_L}^{d-a}\rangle\omega^r\), and
\(G_L=\{u>\psi-L\}\cap\{v>\psi-L\}\). These measures have the same mass
\(M\) and agree on \(G_L\). Their complementary masses therefore agree
as well. On the complement \(B_L\), \(|u_L-v_L|\leq L\). With
\(z_u=\psi-u\), \(z_v=\psi-v\), \[
L\nu_{a,L}(B_L)=L\nu_a(B_L)
\leq\int_{B_L}(z_u+z_v)d\nu_a\longrightarrow0.
\] The last limit follows from mixed weight integrability and the fact
that \(B_L\) decreases to polar sets. The original tail of \(|u-v|\)
also tends to zero. On \(G_L\), the weights agree exactly. Thus every
term in \(\mathcal E_{v_L,\lambda}(u_L)\) converges to the corresponding
finite term in (19a). Taking limits in the bounded cocycle proves (20a).
Additive constants reduce the general case to this normalization.

If \(\psi'\sim\psi\), Vu Corollary 3.2 identifies the finite-energy
classes of currents, and bounded mixed mass equality gives the same
\(d,M\). Representatives of a current differ only by constants on
connected \(X\), so the potential domains agree as well. To check the
energy value without using an unproved truncation interchange, use
\(u_L\) truncated against \(\psi\), and the measure
\(\langle T_u^aT_{\psi'}^{d-a}\rangle\omega^r\). Mixed mass and weight
integrability hold by the preceding finite-family argument, since
\(\psi'\) itself belongs to the domain. On \(\{u>\psi-L\}\) products
agree with those of \(u_L\); the complement has equal mass. Its
truncated weight \(|u_L-\psi'|\) is at most \(L+\|\psi-\psi'\|_\infty\),
and its integral tends to zero by integrability of \(\psi-u\) against
this mixed measure. The original tail also tends to zero. Thus
\(\mathcal E_{\psi',\lambda}(u_L)\to\mathcal E_{\psi',\lambda}(u)\), and
the bounded anchor cocycle passes to the limit: \[
\mathcal E_{\psi',\lambda}(u)=\mathcal E_{\psi,\lambda}(u)-\mathcal E_{\psi,\lambda}(\psi').
\]

Finally let \(\omega'\) be any Kähler representative of the same
\(\lambda\). Closedness and cohomology make each scalar full-mass
condition in (17a) unchanged. Pointwise comparison of the two fixed
Kähler forms makes every absolute-weight mixed integral finite for one
exactly when it is finite for the other. Thus the domain is unchanged.
The bounded energies of every \(u_L\) are representative independent by
n1; (22a) proves representative independence of their finite-energy
limit. A simultaneous smooth-background change as in n1 leaves the
construction unchanged. This proves all reference-change assertions
stated in this note.

\subsection{A domain member with an actually unbounded
difference}\label{a-domain-member-with-an-actually-unbounded-difference}

Take \(X=\mathbb P^1\times\mathbb P^1\), normalized Fubini--Study
\(\kappa\), \(\theta=p_1^*\kappa\), \(\psi=0\), and
\(\omega=p_1^*\kappa+p_2^*\kappa\). Then \(d=1\), \(M=1\). In a disk
about \(0\in\mathbb P^1\), choose \(0<a<1/2\) and \[
v(z)=-b(-\log|z|)^a
\] near zero, with \(b>0\) small. Continue it by a smooth cutoff away
from the disk, and choose \(b\) sufficiently small that
\(\kappa+dd^cv\geq0\); its singular radial Hessian is positive, and the
cutoff contributes a bounded Hessian only on a fixed annulus. Subtract a
constant to make \(v\leq0\). Near zero its positive Laplacian density is
a constant times \(b a(1-a)|z|^{-2}(-\log|z|)^{a-2}\). The boundary flux
is proportional to \((-\log|z|)^{a-1}\to0\), so there is no point atom.
Therefore its nonpluripolar mass is 1. Moreover \[
\int(-v)dd^cv\quad\text{has singular part proportional to}
\quad\int^{\infty}t^{2a-2}dt<\infty,
\] and \(\int(-v)\kappa<\infty\). For \(u=p_1^*v\), both weighted terms
of (17a) are finite and the mass is \(M\), but \(u-\psi\) is unbounded.
This is a genuine finite-energy extension, with no fibration assumed in
the theorem; the product example only demonstrates its domain.

\section{A fixed-class concentration counterexample for unbounded
perturbative
energies}\label{a-fixed-class-concentration-counterexample-for-unbounded-perturbative-energies}

This note uses the bounded coefficient identity of \textbf{n1} and the
finite-energy domain and example of \textbf{n2}, and proves the
additional unbounded application directly. It refutes the proposed
extension of the class-convergent auxiliary-path limit to all
degree-\(d\) finite-energy potentials. It does not refute every possible
unnormalized energy theory in P15.1.

Use \(X=\mathbb P^1\times\mathbb P^1\), \(\theta=p_1^*\kappa\),
\(\psi=0\), and the fixed unbounded \(u=p_1^*v\) constructed in n2, with
\(0<a<1/2\), \(v\leq0\), and \(v(z)=-b(-\log|z|)^a\) near \(0\).
Normalize \(\int\kappa=1\). The class \(\alpha=[\theta]\) is nef,
pseudoeffective, and has volume zero because \(\theta^2=0\) and the
positive product of nef classes equals the cup product. Its actual
reference degree is \(d=1\), and \(u\) has finite degree-one relative
joint energy and full degree-one mass.

For each sufficiently small \(\varepsilon>0\), put \[
\rho_\varepsilon=\exp(-\varepsilon^{-2/a}).
\] Choose a nonnegative smooth two-form \(\beta_\varepsilon\), with
total integral 1, supported in \(\{|z|<\rho_\varepsilon\}\). Such a form
is obtained by rescaling and normalizing any fixed nonnegative smooth
disk bump; it extends smoothly by zero. On a real surface every two-form
is closed. Set \[
\kappa_\varepsilon=\tfrac12\kappa+\tfrac12\beta_\varepsilon,
\qquad
\omega_\varepsilon=p_1^*\kappa_\varepsilon+p_2^*\kappa.
\tag{23a}
\] These are smooth Kähler forms. Each \([\kappa_\varepsilon]=[\kappa]\)
on \(\mathbb P^1\), because top real cohomology is determined by the
integral. Thus all \([\omega_\varepsilon]\) equal the single Kähler
class \(\lambda=[p_1^*\kappa+p_2^*\kappa]\), although their
representatives do not converge smoothly.

Write \(S=\kappa+dd^cv\) on the base. It is positive of mass 1, with no
polar atom. The current \(T_u=p_1^*S\) is positive; thus the fixed
choice \(u_\varepsilon=u\) is
\((\theta+\varepsilon\omega_\varepsilon)\)-psh and is uniformly
identical to \(u\). For each fixed \(\varepsilon\), all relevant
unbounded weighted integrals are finite: \(v\) is integrable against the
smooth \(\kappa_\varepsilon\), and against both \(S\) and \(\kappa\) by
n2. Each has a constant depending on that fixed form, which need not be
uniform in \(\varepsilon\). Locality off \(z=0\) and the absence of
polar mass give \[
\langle(T_u+\varepsilon\omega_\varepsilon)^2\rangle
=2\varepsilon\,p_1^*(S+\varepsilon\kappa_\varepsilon)\wedge p_2^*\kappa.
\] Its total mass is \(2\varepsilon(1+\varepsilon)\), exactly the
anchor's mass. Hence \(u\) is of full relative top mass for every
positive perturbation, and its relative top joint linear-weight energy
is finite for every \(\varepsilon>0\). This counterexample does not
exploit lost positive-perturbation mass or infinite energy at a fixed
parameter.

Directly expand the three mixed terms in the unnormalized top energy
with anchor zero: \[
\begin{aligned}
(T_0+\varepsilon\omega_\varepsilon)^2
 &=2\varepsilon p_1^*(\kappa+\varepsilon\kappa_\varepsilon)p_2^*\kappa,\\
\langle(T_u+\varepsilon\omega_\varepsilon)(T_0+\varepsilon\omega_\varepsilon)\rangle
 &=\varepsilon p_1^*(S+\kappa+2\varepsilon\kappa_\varepsilon)p_2^*\kappa,\\
\langle(T_u+\varepsilon\omega_\varepsilon)^2\rangle
 &=2\varepsilon p_1^*(S+\varepsilon\kappa_\varepsilon)p_2^*\kappa.
\end{aligned}
\] All weighted terms are finite, so expansion under their integrals is
legitimate. Consequently \[
\mathcal E^{2,\varepsilon,\omega_\varepsilon}_0(u)
=2\varepsilon\mathcal E_{0,\lambda}(u)
 +2\varepsilon^2\int_{\mathbb P^1}v\kappa_\varepsilon,
\qquad
\mathcal E_{0,\lambda}(u)=\tfrac12\int v(S+\kappa)\in\mathbb R.
\tag{24a}
\] Thus the normalized relative energy is \[
\frac{\mathcal E^{2,\varepsilon,\omega_\varepsilon}_0(u)}{2\varepsilon}
=\mathcal E_{0,\lambda}(u)+\varepsilon\int v\kappa_\varepsilon.
\tag{25a}
\] On the support of \(\beta_\varepsilon\), one has
\(-v\geq b\varepsilon^{-2}\), possibly strengthened by the harmless
constant used to normalize \(v\). Therefore \[
\int v\kappa_\varepsilon\leq-\tfrac b2\varepsilon^{-2},
\qquad
\varepsilon\int v\kappa_\varepsilon\longrightarrow-\infty.
\] The expression (25a) diverges to minus infinity. For the fixed smooth
auxiliary form \(\omega=p_1^*\kappa+p_2^*\kappa\), the same computation
has the fixed finite \(\int v\kappa\), and the expression converges to
\(\mathcal E_{0,\lambda}(u)\). Both paths have exactly the same
auxiliary cohomology class and exactly the same potential for every
parameter.

This is an analytic counterexample, within all ambient hypotheses, to an
\textbf{unbounded strengthening} of n1's class-convergent perturbation
theorem. It is not a counterexample to n1, whose bounded-difference
hypothesis supplies (14). It shows that degree-\(d\) finite energy, even
together with finite full relative top energy at every positive
perturbation, does not provide the uniform weighted lower-coefficient
estimate needed for an auxiliary representative path. Here the lower
weighted coefficient is finite at every parameter and still diverges
faster than its suppressing power of \(\varepsilon\). Additional uniform
integrability, a restriction on auxiliary forms, or a different
iterated/truncated limiting prescription is necessary for such a path
theorem. None is silently added to the original catalogue atom.

\section{Reference data, exact residual scope, and execution
evidence}\label{reference-data-exact-residual-scope-and-execution-evidence}

This note uses local \textbf{n1}, \textbf{n2}, and \textbf{n3}. It gives
a bounded-family characterization of the necessary reference scale,
explicit examples, and the remaining scope of P15.1. It does not claim
original closure.

\subsection{A proved characterization of scale equivalence on the
bounded
domain}\label{a-proved-characterization-of-scale-equivalence-on-the-bounded-domain}

Fix the actual bounded singularity type \(\tau=[\psi]\), and its degree
\(d\); let \(r=n-d\). Put \(V=H^{1,1}(X,\mathbb R)\). For
\(u,v\in\mathcal C_\psi\), define a symmetric multilinear functional on
\(r\) Kähler classes by \[
L_{u,v}(\lambda_1,\ldots,\lambda_r)
=\frac1{d+1}\sum_{a=0}^d\int(u-v)
\langle T_u^aT_v^{d-a}\rangle\omega_1\cdots\omega_r,
\quad\omega_i\in\lambda_i.
\tag{26a}
\] It is independent of each representative: changing one factor inserts
\(dd^cg\), the other \(r-1\) factors are closed positive, and the
admissible IBP and degree-\(d+1\) vanishing argument (11) applies
verbatim. Extending factors as differences of Kähler forms gives a
well-defined symmetric multilinear functional on all of \(V^r\). To
justify well-definedness of this extension explicitly, represent a real
class \(a\) by a smooth closed real form \(\eta\). Both \(\eta+C\Omega\)
and \(C\Omega\) are Kähler for sufficiently large \(C\). The defining
integral is multilinear in smooth factors; any two decompositions have
the same difference form up to \(dd^c\) of a smooth function.
Representative invariance for positive factors, then expansion by
differences in all remaining factors, makes that change zero. Thus there
is a linear functional \(\ell_{u,v}\) on \(\operatorname{Sym}^rV\) with
value (26a). Constants in the potential are included among \(u,v\).

Set \[
K_\tau=\bigcap_{u,v\in\mathcal C_\psi}\ker\ell_{u,v}.
\tag{27a}
\] Then the bounded energy family needs only the image of the formal
tensor \(\lambda^{\odot r}\) in the finite-dimensional quotient
\(\operatorname{Sym}^rV/K_\tau\). Conversely two formal tensors give the
same values for \textbf{every} bounded relative energy exactly when
their difference belongs to \(K_\tau\), by definition of the
intersection of kernels. This is a proved coarsest scale-equivalence
characterization for this fixed bounded family, not a declaration that
the entire auxiliary class is always minimal. Changing the anchor within
\(\tau\) does not change \(K_\tau\), since all relative differences
remain in the same family. The degree-zero case uses the same formula
with the empty current product and \(r=n\).

The characterization does not construct a geometric lower-dimensional
carrier, nor a universal minimal scale for all singularity types or all
unbounded domains. For the finite-energy extension n2, a fixed Kähler
class is sufficient and its representatives are unnecessary, but no
minimal quotient characterization across arbitrary classes is claimed.

\subsection{Why a scale and actual type cannot be
discarded}\label{why-a-scale-and-actual-type-cannot-be-discarded}

On \(X=\mathbb P^1\times\mathbb P^1\), with normalized \(\kappa\), let
\(\alpha=[p_1^*\kappa]\), \(\theta=p_1^*\kappa\), and \(\psi=0\). It is
a nef zero-volume class by the nef cup-product statement. For
\(\omega=A p_1^*\kappa+B p_2^*\kappa\), \(A,B>0\), one has \(d=1\),
\(M_1=B\), and \(\mathcal E_{0,[\omega]}(c)=cB\). Varying \(B\) changes
the raw energy of the same constant potential. Some scale is therefore
necessary. Varying \(A\) alone does not change these energies; it
illustrates why the entire Kähler class should not simply be declared
minimal. Indeed all bounded \(\theta\)-psh potentials descend to the
first factor: their restriction to each second-factor fiber is a psh
function on a compact curve and is constant, and bounded potentials are
finite on every fiber. Equivalently the constant-shift example alone
already proves necessity of \(B\), without any descent theorem.

Actual reference singularity data also affect the degree. On the first
\(\mathbb P^1\), let \(g(z)=\log(|z|^2/(1+|z|^2))\), in the chart whose
distinguished point is zero. In the chosen convention
\(\kappa+dd^cg=\delta_0\): the elementary Poincaré--Lelong identity
\(dd^c\log|z|^2=\delta_0\) follows by the distributional boundary flux
around a disk, and \(dd^c\log(1+|z|^2)=\kappa\). Thus the actual
reference \(\psi=p_1^*g\) represents the positive integration current of
the vertical divisor. Its degree-one nonpluripolar current is zero,
since it is supported entirely on its polar divisor; all higher products
vanish as well. This reference has \(d=0\) in the same class whose
smooth reference has \(d=1\). The numerical dimension of this nef class
is 1 by its nonzero first cup product and zero square, but the actual
prescribed reference can have degree zero. Hence the degree in n1 is not
identified with the class numerical dimension.

In this same example the actual perturbation mass is
\(P_\varepsilon=2AB\varepsilon^2\), whereas the class is nef and its
perturbation volume is the cup product
\(\operatorname{vol}(\alpha+\varepsilon[\omega])=2B\varepsilon+2AB\varepsilon^2\).
The first nonzero class-volume coefficient therefore has a different
exponent from the actual-reference mass. An actual anchor constant shift
has relative perturbed energy \(cP_\varepsilon\); rescaling by the
class-volume leading exponent would make its limit zero. This is an
explicit reason not to substitute a class-volume coefficient for the
actual prescribed-type scale without proving an additional theorem.

The actual carrier throughout is \(X\), of complex dimension \(n\); the
integrals use a degree-\(d\) nonpluripolar current and complementary
smooth factors. The full auxiliary form would remain data for its
differential measure. The energy itself only uses the class, or the
precise quotient (27a) within the bounded family. An actual anchor fixes
the additive origin of energy; changing it has the stated cocycle. Gauge
changes of the smooth background and potentials leave every current
unchanged.

\subsection{What remains open for the original
P15.1}\label{what-remains-open-for-the-original-p15.1}

The bounded theorem removes the supplied-fibration hypothesis
completely: its ambient class and actual reference are arbitrary as in
the question. The finite-energy extension is a genuine more-singular
domain, not just bounded differences under another name. Nevertheless
these are scoped partial results; they are not counted as original
catalogue closure.

The unbounded extension of the class-convergent auxiliary-path
prescription is \textbf{false} without further conditions, by the
analytic counterexample n3. Even within the finite-energy domain, lower
weighted perturbation coefficients can lack uniform control along
fixed-class smooth representatives. For general references and fixed
auxiliary forms, integrability and uniform control of all lower-degree
weighted terms, and relative limits for genuinely unbounded moving
potentials, have not been proved here. Nor have continuity for arbitrary
increasing sequences, differentiation with unbounded velocities, or all
proposed path classes been established. Decreasing continuity, cocycles,
fixed-class representative invariance, and concavity in n2 are the
precise asserted extension properties.

There is no numerical-dimension theorem, no general leading
\textbf{class-volume} coefficient theorem, no complete universal
minimum-reference construction for unbounded types, and no intrinsic
geometric \(d\)-dimensional carrier. The degree and perturbation
exponent use actual-current masses rather than an assumed equality with
class volume. No fibration, ideal-equivalence-only recovery, zero-mass
probability normalization, nonarchimedean translation, or independent
worker verdict is used. P15.2 and P15.3 are outside this assignment. The
original P15.1 remains open in the dossier's residual scope.

\subsection{Audit and rerun contract}\label{audit-and-rerun-contract}

The original \texttt{input.json} is retained byte-for-byte, with SHA256
\path{5950ca2b12dc2202093de059446f0076db58a9cc861e1e57ce8c82edc371753b}.
Its decoded assignment string is also retained as
\texttt{assignment-original.txt}; \texttt{task-original.json} preserves
the task values. The authoritative deliverable is
\texttt{P15.1-argument.md}, containing all four mathematical notes and
the complete immutable source bodies. \texttt{P15.1-partial.tex} and its
twice-compiled \texttt{P15.1-partial.pdf} contain the entire new proof,
standard foundation statements and applicability checks. Literal source
bodies remain in the Markdown annex to preserve their UTF-8 extraction
and defects.

The independent source audit reads only
\texttt{/workspace/research/shared-sources}. All 14 supplied LF-range
SHA256 pins match exactly when a range is joined by literal LF without
appending an extra final LF. The audit additionally records three
directly read foundation ranges. All source files are read-only and
unchanged. Source text retains form feeds, the printed missing \(1/2\),
the v1 Proposition 4.2 \(m/n\) and subscript defects, and the malformed
extracted convention. The correct intrinsic formula and all new
calculations are separate from these raw quotations. The page-4/page-6
Vu v3 images and page-4 Vu v1 convention image were inspected locally.
Vu v1's finite-energy theorem is not the controlling finite-energy
source. No source-dependent catalogue claim was treated as independently
established by the catalogue description.

The requested execution identity was native Codex \texttt{gpt-6.1-sol},
reasoning \texttt{xhigh}. This turn used the native available Codex
runtime only; no additional model or subagent was invoked. The runtime
does not expose a signed model/effort attestation here, so those exact
backend identifiers are recorded as requested rather than independently
verified. No network call, install, external message, auth/proxy/CA
change, HOME/CODEX\_HOME change, SQL/runtime edit, or old-artifact write
was performed. All new files are in this assigned workspace. There was
no resource-heavy computation and no need to invoke a supervised compute
service.

Rerun commands from this directory are:

\begin{itemize}
\tightlist
\item
  \texttt{python\ verify\_sources.py}
\item
  \texttt{python\ coefficient\_check.py}
\item
  \texttt{python\ build.py}
\end{itemize}

Observed checks report 14/14 matching supplied source pins, 968 passing
exact finite Vandermonde coefficient cases, and two successful
\texttt{pdflatex} passes followed by successful PDF text-extraction
checks. The initial compilation encountered an authored inline-math
delimiter error; \texttt{repair\_math\_markup.py} corrected the
delimiters without changing mathematical content, and
\texttt{pdflatex-initial-failed.log} preserves that diagnostic. The
final document was then compiled successfully twice. The build records
environment versions, commands, return codes, source hashes,
twice-compiled PDF evidence, and diagnostic output in local JSON/log
artifacts. Coefficient checks are exact finite arithmetic diagnostics,
not analytic proofs. PDF compilation and artifact validation are
observations, not mathematical correctness certification.
\texttt{status.json} has exactly the five requested keys and records
\texttt{status:\ partial}. \texttt{notes.json} is the canonical native
note object with \texttt{candidate:\ false}; its support arrays contain
only earlier local IDs. No literature-priority claim or external
independent correctness verdict is made.

\end{document}
